\section{Idempotents and Module Homomorphisms}

\begin{theorem}
    Let $R$ be a ring, let $e,f\in R$ be idempotents, and let ${}_RV$ be a left $R$-module.
    \begin{enumerate}
        \item There is a natural isomorphism of abelian groups
              \[
                  \Hom_R(Re,V)\cong eV.
              \]
              In particular,
              \[
                  \Hom_R(Re,Rf)\cong eRf.
              \]

        \item There is an isomorphism of rings
              \[
                  \operatorname{End}_R(Re)\cong (eRe)^{op}.
              \]
    \end{enumerate}
\end{theorem}

\begin{proof}
    \textbf{(1)} Define
    \[
        g:\Hom_R(Re,V)\longrightarrow eV,\qquad \varphi\longmapsto \varphi(e).
    \]
    First we check that $g$ is well-defined. If $\varphi\in \Hom_R(Re,V)$, then
    \[
        \varphi(e)=\varphi(e^2)=e\varphi(e),
    \]
    so indeed $\varphi(e)\in eV$.

    Next, for any $a\in R$ and any $\varphi\in \Hom_R(Re,V)$, we have
    \[
        \varphi(ae)=a\varphi(e).
    \]
    Thus $\varphi$ is uniquely determined by the single value $\varphi(e)$. Hence $g$ is injective.

    To show surjectivity, let $v\in eV$. Define
    \[
        \psi_v:Re\longrightarrow V,\qquad ae\longmapsto av.
    \]
    This is well-defined, since if $ae=be$, then
    \[
        av=(ae)v=(be)v=bv.
    \]
    It is also $R$-linear. Moreover,
    \[
        g(\psi_v)=\psi_v(e)=ev=v,
    \]
    because $v\in eV$. Hence $g$ is surjective.

    Finally, $g$ is a homomorphism of abelian groups, since for $\varphi,\psi\in \Hom_R(Re,V)$,
    \[
        g(\varphi+\psi)=(\varphi+\psi)(e)=\varphi(e)+\psi(e)=g(\varphi)+g(\psi).
    \]
    Therefore $g$ is an isomorphism of abelian groups:
    \[
        \Hom_R(Re,V)\cong eV.
    \]

    Taking $V=Rf$, we obtain
    \[
        \Hom_R(Re,Rf)\cong e(Rf)=eRf.
    \]

    \medskip

    \textbf{(2)} Apply part~\textbf{(1)} with $V=Re$. Then we get an isomorphism of abelian groups
    \[
        \Phi:\operatorname{End}_R(Re)\longrightarrow eRe,\qquad \sigma\longmapsto \sigma(e).
    \]
    Now let $\sigma,\tau\in \operatorname{End}_R(Re)$. Then
    \[
        \Phi(\sigma\circ\tau)
        =(\sigma\circ\tau)(e)
        =\sigma(\tau(e)).
    \]
    Since $\tau(e)\in Re$, we have $\tau(e)=\tau(e)e$, and hence
    \[
        \sigma(\tau(e))
        =\sigma(\tau(e)e)
        =\tau(e)\sigma(e).
    \]
    Therefore
    \[
        \Phi(\sigma\circ\tau)=\Phi(\tau)\Phi(\sigma).
    \]
    So composition in $\operatorname{End}_R(Re)$ corresponds to multiplication in $eRe$ with the order reversed. Hence
    \[
        \operatorname{End}_R(Re)\cong (eRe)^{op}
    \]
    as rings.
\end{proof}

\begin{theorem}
    Let $e,f$ be idempotents of a ring $R$. Then the followings are equivalent:
    \begin{enumerate}
        \item $Re \cong Rf$ as left $R$-modules.
        \item $eR \cong fR$ as right $R$-modules.
        \item There exist elements
              \[
                  a \in eRf, \qquad b \in fRe
              \]
              such that
              \[
                  ab=e, \qquad ba=f.
              \]
    \end{enumerate}
\end{theorem}

\begin{proof}
    By passing to the opposite ring $R^{\operatorname{op}}$, statement \textup{(2)} is the right-handed analogue of \textup{(1)}. Thus it suffices to prove that \textup{(1)} and \textup{(3)} are equivalent.

    \medskip

    \noindent
    \textup{(1) $\Rightarrow$ (3).}
    Assume that there is an isomorphism of left $R$-modules
    \[
        \varphi: Re \xrightarrow{\;\sim\;} Rf.
    \]
    Set
    \[
        a:=\varphi(e), \qquad b:=\varphi^{-1}(f).
    \]
    Then
    \[
        a=\varphi(e)=\varphi(ee)=e\varphi(e)=ea,
    \]
    and since $a\in Rf$, we also have $af=a$. Hence
    \[
        a\in eRf.
    \]
    Similarly,
    \[
        b=\varphi^{-1}(f)=\varphi^{-1}(ff)=f\varphi^{-1}(f)=fb,
    \]
    and since $b\in Re$, we have $be=b$. Therefore
    \[
        b\in fRe.
    \]

    Now compute:
    \[
        f=\varphi\bigl(\varphi^{-1}(f)\bigr)=\varphi(b)=\varphi(be)=b\varphi(e)=ba,
    \]
    and
    \[
        e=\varphi^{-1}(\varphi(e))=\varphi^{-1}(a)=\varphi^{-1}(af)=a\varphi^{-1}(f)=ab.
    \]
    Thus
    \[
        ab=e, \qquad ba=f.
    \]

    \medskip

    \noindent
    \textup{(3) $\Rightarrow$ (1).}
    Assume there exist $a\in eRf$ and $b\in fRe$ such that
    \[
        ab=e, \qquad ba=f.
    \]
    Since $a\in eRf$ and $b\in fRe$, we have
    \[
        ea=a, \qquad af=a, \qquad fb=b, \qquad be=b.
    \]

    Define maps
    \[
        \varphi_a: Re \longrightarrow Rf, \qquad xe \longmapsto xa,
    \]
    and
    \[
        \varphi_b: Rf \longrightarrow Re, \qquad xf \longmapsto xb.
    \]
    These are well-defined left $R$-module homomorphisms.

    Now for any $xe\in Re$,
    \[
        (\varphi_b\circ \varphi_a)(xe)=\varphi_b(xa)=xab=xe,
    \]
    since $ab=e$. Likewise, for any $xf\in Rf$,
    \[
        (\varphi_a\circ \varphi_b)(xf)=\varphi_a(xb)=xba=xf,
    \]
    since $ba=f$.

    Hence $\varphi_a$ and $\varphi_b$ are inverse isomorphisms, so
    \[
        Re \cong Rf
    \]
    as left $R$-modules.

    Therefore \textup{(1)} and \textup{(3)} are equivalent, and by symmetry so are all three statements.
\end{proof}

